\section{Conclusion}
\label{sec:conclusion}

In this paper, we introduced Q-FUSION, an adaptive classical-quantum hybrid architecture designed to operate under strict resource and latent-dimension constraints in 3D medical vision pipelines. 
By resolving classical collapse through annealed routing, bypassing the spatial reconstruction paradox via quantum-conditioned FiLM modulation, and enforcing strict 0.0000\% capacity parity against classical surrogates, we established an empirically grounded evaluation of variational representations in neuro-oncology. 
Evaluations across 150 BraTS2021 subjects demonstrated that adaptive routing successfully overcomes the multi-modality saturation and feature dilution observed in classical baselines at $M=4$ sequences (+4.79\% DSC advantage). 
On SynthRAD2023, adaptive routing established a +2.47\% mean DSC advantage over rigid quantum baselines across multiple random seeds, while demonstrating flat zero-shot degradation under open-system gate depolarizing noise. 
These results confirm that dynamic classical-quantum representation allocation provides measurable utility and architectural resilience in resource-constrained multimodal medical image analysis.
